\section{The MIDI Protocol}\label{sec:background-midi}

The Musical Instrument Digital Interface (MIDI) is an industry-standard technical protocol developed in 1983
to enable electronic instruments, computers, and other related devices to communicate and synchronize with each other.
Crucially, MIDI does not transmit digital audio (sound waves); instead, it transmits \textit{event messages}
that act as instructions for how music should be produced by a receiving synthesizer.

\subsection{Event-Based Representation}

A MIDI stream is composed of sequential commands. The most fundamental commands are:
\begin{itemize}
  \item \textbf{Note-On:} Instructs the synthesizer to begin sounding a pitch. It carries two data bytes: the
    \textit{key number} (0--127, where 60 is typically Middle C) and the \textit{velocity} (0--127,
    representing how hard the key was struck).
  \item \textbf{Note-Off:} Instructs the synthesizer to stop sounding the pitch.
  \item \textbf{Program Change:} Changes the active instrument (or "patch") on a specific channel (e.g.,
    switching from a piano to a violin).
\end{itemize}

Because MIDI transmits events rather than waveforms, a piece of music represented in MIDI is highly editable.
A composer can retroactively change the pitch, duration, or instrument of any note without degrading the
audio quality, making MIDI the \textit{de facto} standard for digital composition.

\subsection{Channels and General MIDI}

To support polyphony across different instruments, MIDI data is transmitted across 16 independent channels.
A synthesizer can be configured to play a bass patch on Channel 1 and a piano patch on Channel 2 simultaneously.
The \textit{General MIDI} (GM) specification standardizes the assignment of instruments to program numbers
(e.g., Program 1 is Acoustic Grand Piano) and strictly reserves \textbf{Channel 10} for unpitched percussion,
where individual key numbers trigger different drum samples (e.g., Key 36 is the Bass Drum).

\subsection{Delta-Times and Standard MIDI Files (SMF)}

When saving MIDI data to a file (a \texttt{.mid} file), the protocol utilizes \textit{delta-times}.
Instead of recording the absolute time at which an event occurs, a Standard MIDI File records the number of
"ticks" to wait after the \textit{previous} event before executing the current one.
Ticks are subdivisions of a quarter note, defined by the file's \textit{Pulses Per Quarter-note} (PPQ)
resolution (e.g., 480 PPQ).
This relative timing system is highly efficient but requires careful calculation when translating from
absolute musical beats into a serialized file format.
